% Generated draft: requires mathematical and source review.
\documentclass[11pt]{article}
\usepackage[margin=25mm]{geometry}
\usepackage{fontspec}
\setmainfont{Latin Modern Roman}
\usepackage{amsmath,amssymb,mathtools,unicode-math}
\setmathfont{Latin Modern Math}
\usepackage{enumitem,needspace,xurl,hyperref}
\setlength{\parindent}{0pt}
\setlength{\parskip}{5pt}
\setlength{\emergencystretch}{3em}
\newcommand{\N}{\mathbb{N}}
\newcommand{\Z}{\mathbb{Z}}
\newcommand{\Q}{\mathbb{Q}}
\newcommand{\R}{\mathbb{R}}
\newcommand{\C}{\mathbb{C}}
\newcommand{\F}{\mathbb{F}}
\newcommand{\A}{\mathbb{A}}
\newcommand{\PP}{\mathbb{P}}
\newcommand{\E}{\mathbb{E}}
\newcommand{\eps}{\varepsilon}
\newcommand{\dd}{\,\mathrm{d}}
\begin{document}
\section*{UnsolvedMath ID: 20001983}
\textbf{Source problem code:} catalogue entry
\subsection{Definitions and mathematical statement}
Federico Ardila's question asks for combinatorial descriptions of Bergman complexes beyond linear subspaces, analogous to the matroid description of Ardila--Klivans. For a matroid \(\mathcal M\), let \(L(\mathcal M)\) be its lattice of flats. The relevant order complex uses the proper part, with the bottom and top flats removed: its simplices are chains of proper nonempty flats. In the usual loopless linear-space setting, the source reports that this gives a subdivision of the Bergman complex, and hence access to its topology. The excerpt calls a general variety's tropical variety its Bergman complex; fan/link, lineality, and valuation conventions are not specified. The broad request is to find other classes with comparable descriptions, not a single quantified universal conjecture.

The aid reports a combinatorial answer for tropical hypersurfaces. To state its standard polyhedral model precisely, let \(f=\sum_{a\in A}c_ax^a\) be a Laurent polynomial over a nonarchimedean valued field, with finite support \(A\subset\Z^n\), and assume \(P=\operatorname{conv}(A)\) has dimension \(n\). Use the minimum convention
\[
 \operatorname{Trop}(V(f))=\{w\in\R^n:
 \min_{a\in A}(\operatorname{val}(c_a)+\langle a,w\rangle)
 \text{ is attained at least twice}\}.
\]
The coefficient heights induce a regular subdivision of \(P\). The reported duality reverses face inclusion between its positive-dimensional faces \(\sigma\) and tropical cells \(\sigma^\vee\), with
\[
 \dim\sigma+\dim\sigma^\vee=n.
\]
The cell \(\sigma^\vee\) is bounded exactly when \(\sigma\) is not contained in \(\partial P\). Full dimension is explicit here; lower-dimensional Newton polytopes introduce lineality and are not covered by this boundedness formulation without adjustment.

The aid's worked family takes integers \(n\ge2\), \(d>n\), an element \(t\) with \(\operatorname{val}(t)=1\), and
\[
 f_{n,d}=t(1+x_1^d+\cdots+x_n^d)+x_1\cdots x_n.
\]
Its Newton polytope is \(d\Delta_n=\operatorname{conv}(0,de_1,\ldots,de_n)\), with the additional support point \(q=(1,\ldots,1)\) strictly interior. The aid reports that the coefficient valuations induce the star triangulation at \(q\). Consequently the bounded tropical face poset is opposite to the nonempty face poset of the boundary of an \(n\)-simplex. Since that simplex boundary is combinatorially self-dual, the bounded complex is reported to be combinatorially a simplex boundary, with
\[
 |\operatorname{Trop}(V(f_{n,d}))_{\mathrm{bounded}}|\cong S^{n-1},\qquad
 f_k=\binom{n+1}{k+1}\quad(0\le k\le n-1).
\]
Here \(f_k\) counts bounded \(k\)-cells, and only the bounded complex is asserted to be this sphere. The aid singles out \((n,d)=(2,3)\) for an audit of every cell, weight, and balancing relation, but those audit details are not supplied. This is a nonlinear valued-hypersurface answer to part of the broad request, not an identification of every tropical variety with a matroid order complex.
\subsection{Short English statement}
Ardila asks for combinatorial models of tropical varieties beyond the matroid model for linear spaces. The aid uses regular Newton subdivisions for hypersurfaces and reports a sparse family whose bounded tropical complex is a simplex boundary. Valuation, link, lineality, and bounded-complex conventions remain distinct from the original loosely specified Bergman-complex terminology.
\subsection{Sources}
\begin{itemize}
\item Local UnsolvedMath registry, permanent ID 20001983, exact source code catalogue entry; canonical \texttt{aim-geometry-notes.json}, \texttt{notes[320]}, tag \texttt{section}; AIM Workshop Problem Lists, Geometry, workshop \emph{Amoebas and tropical geometry}, A.1, contributed by Federico Ardila. Created and updated 14 August 2026.
\item AIM source \url{https://aimath.org/WWN/amoebas/amoebas.pdf}, page 3 according to supplied verification; it reports agreement with an AIM HTML transcription, without supplying a separate HTML URL. Neither external source was accessed.
\item Sturmfels, \emph{Solving systems of polynomial equations}, Chapter 9, cited for the Bergman/tropical notion; a paper of Ardila and Klivans, cited for the lattice-of-flats subdivision. No complete publication entry for that paper or precise theorem reference for Newton duality is supplied.
\item Supplied attempt 1: \texttt{partial\_result}, problem status at run \texttt{partially\_solved}, rigor label \texttt{proved}, novelty confidence \texttt{low}; classification \texttt{PARTIAL-PROGRESS}, registry status \texttt{partially\_solved}. Candidate: sparse all-dimensional worked family and face enumeration. The claim that the hypersurface class has a complete known answer is the aid's supplied assessment.
\item Ulam AI / UnsolvedMath Contributors; CC BY 4.0. Machine-generated research aid, not peer reviewed; independently verify before citation.
\end{itemize}
\end{document}
